% page 549 du fac-simile (image p-548 du scan)
% bloc 170.2 x 242.0 mm ; marge gauche 31.8 mm
\pgc{10.160mm}{0.000mm}{
\l{\cel{55}541}
\l{}
\l{\sou{solicitar}. (\sou{trans., pri}). I. Incitar (ulo) tre insistoze pri ulo.}
\l{\cel{3}- II. Pregar (ulu) tre insistoze por obtenar (de lu) ulo. -}
\l{\cel{3}DEFIS.}
\l{}
\l{\sou{solida}. Qua konsistas ye molekuli inter-agregita ye maniero kohe-}
\l{\cel{3}ranta, ferma. - DEFIS.}
\l{}
\l{\sou{solidara}. I. (\sou{yurocienco}). Tale komuna ye pluri ke singla respon-}
\l{\cel{3}sas pri la ensemblo. - II. Ligita a kunhomi per responso komuna.}
\l{\cel{3}- DEFIRS.}
\l{}
\l{\sou{solipeda}. (\sou{zool.)}\cel{2}Qua havas nur un fingro hufoza ye singla ek sua}
\l{\cel{3}pedi. - DEFIS.}
\l{}
\l{\sou{solitara}. Loko solitara : ube onu esas sola. - EFIS.}
\l{}
\l{\sou{solitero}. I. Diamanto qua esas inkastrita sole. - II. Ludilo ek}
\l{\cel{3}tabuleto qua kontenas 37 trui, ye qua onu ludas (sole) per 36}
\l{\cel{3}stifti. - DEFIRS.}
\l{}
\l{\sou{solstico}.\cel{2}(\sou{astron.)}\cel{2}Singla ek la du epoki en qui Suno distas}
\l{\cel{3}maxime de equatoro e semblas sen-mova dum kelka dii. - DEFIRS.}
\l{}
\l{\sou{solvar}. (\sou{trans.})\cel{2}I. Deskompozar e retro-igar ad stando elementala.}
\l{\cel{3}II. Desintrikar de to quo embarasas la spirito. - EFIS.}
\l{}
\l{\sou{solventa}. (\sou{komerco)}. Qua povas pagar sua debaji. - DEFIRS.}
\l{}
\l{\sou{somato}. (\sou{biol.)}\cel{2}Ensemblo ek la celuli qui desaparas e mortas}
\l{\cel{3}samatempe kam la individuo (kontraste kun la celuli jerma qui}
\l{\cel{3}persistas sen-fine per genito). - DEFIS.}
\l{}
\l{\sou{somero}. Sezono kun vetero varma, qua sequas printempo e preiras}
\l{\cel{3}autuno. - DE.}
\l{}
\l{\sou{somito}. La parto maxim supra di to quo stacas sur sua bazo. - EFI.}
\l{}
\l{\sou{somnambula}. Dicesas pri homo qua, en la stando ne-normala quan}
\l{\cel{3}produktas lua tendenco naturala o la efiko di magnetizo, agas}
\l{\cel{3}e parolas dum sua dormo. - DEFIS.}
\l{}
\l{\sou{somnolar}. (\sou{netrans.}) Sentar sua bezono quik dormor. - DEFIS.}
\l{}
\l{\sou{sonar}. (\sou{netrans.}) (\sou{aludante objekto o materio)}.\cel{2}Emisar movi}
\l{\cel{3}vibrala, propagate da medio aeroforma, liquida o solida, til}
\l{\cel{3}la membrano di la orel-timpano, produktas sento partikulara.}
\l{\cel{3}- EFIS.}
\l{}
\l{\sou{sonalio}. Mikra bulo, kava, ek metalo, qua kontenas frapilo moviva}
\l{\cel{3}qua igas lu resonigar lor movo mem mikra. - I.}
\l{}
\l{\sou{sonato}.- (\sou{muziko}).Muzikajo klasika, por un o plura instrumenti,}
\l{\cel{3}ordinare ek \textquotedbl{}prefaco\textquotedbl{}, \textquotedbl{}adagio\textquotedbl{} o \textquotedbl{}andante\textquotedbl{}, menueto o finalo.}
\l{\cel{3}- DEFIRS.}
}
